\documentclass[11pt,letterpaper]{article}
\usepackage[T1]{fontenc}
\usepackage[utf8]{inputenc}
\usepackage{amsmath,amsthm,mathtools}
\usepackage{newtxtext,newtxmath}
\usepackage[margin=1in,headheight=15pt]{geometry}
\usepackage{microtype,booktabs,array,longtable,enumitem,xurl}
\usepackage{xcolor,fancyhdr,needspace}
\usepackage[colorlinks=true,linkcolor=ink,citecolor=ink,urlcolor=ink]{hyperref}
\usepackage[nameinlink,noabbrev]{cleveref}
\definecolor{ink}{HTML}{354E42}
\definecolor{pale}{HTML}{F1F2E9}
\hypersetup{pdftitle={Forcing, Omitted Cuts, and Old Surreal Fields},
 pdfauthor={Research article prepared with ChatGPT},
 pdfsubject={External saturation, ordinal sequences, surcomplex conjugation, and Boolean completions}}
\pagestyle{fancy}\fancyhf{}
\fancyhead[L]{\small\scshape Forcing and old surreal fields}
\fancyhead[R]{\small Set theory and model theory}
\fancyfoot[C]{\thepage}
\renewcommand{\headrulewidth}{0.3pt}
\setlength{\parskip}{3.5pt}
\setlength{\parindent}{1.15em}
\setlength{\emergencystretch}{2em}
\setlist{leftmargin=1.7em,itemsep=2pt,topsep=5pt}
\clubpenalty=10000\widowpenalty=10000
\allowdisplaybreaks[2]
\numberwithin{equation}{section}
\newtheorem{theorem}{Theorem}[section]
\newtheorem{lemma}[theorem]{Lemma}
\newtheorem{proposition}[theorem]{Proposition}
\newtheorem{corollary}[theorem]{Corollary}
\theoremstyle{definition}
\newtheorem{definition}[theorem]{Definition}
\newtheorem{example}[theorem]{Example}
\theoremstyle{remark}
\newtheorem{remark}[theorem]{Remark}
\newtheorem{question}[theorem]{Question}
\newcommand{\No}{\mathbf{No}}
\newcommand{\Ord}{\mathrm{Ord}}
\newcommand{\R}{\mathbb R}
\newcommand{\C}{\mathbb C}
\newcommand{\Q}{\mathbb Q}
\newcommand{\N}{\mathbb N}
\newcommand{\Z}{\mathbb Z}
\newcommand{\Pow}{\mathcal P}
\newcommand{\bd}{\operatorname{b}}
\newcommand{\cf}{\operatorname{cf}}
\newcommand{\acl}{\operatorname{acl}}
\newcommand{\rcl}{\operatorname{rcl}}
\newcommand{\dom}{\operatorname{dom}}
\newcommand{\ran}{\operatorname{ran}}
\newcommand{\res}{\mathbin{\upharpoonright}}
\newcommand{\NS}{\mathsf{NS}}
\newcommand{\RCF}{\mathsf{RCF}}
\newcommand{\ACF}{\mathsf{ACF}_0}
\newcommand{\Lang}{\mathcal L}
\newcommand{\abs}[1]{\lvert#1\rvert}
\newcommand{\cut}[2]{\left\{\,#1\;\middle|\;#2\,\right\}}
\newcommand{\seq}[1]{\langle#1\rangle}
\newcommand{\shorthash}{\texttt{4f2645599121}}
\newcommand{\fullhash}{\texttt{4f264559\allowbreak 9121fa87\allowbreak 2c710499\allowbreak 5e47f86f\allowbreak 0382351f}}
\newcolumntype{P}[1]{>{\raggedright\arraybackslash}p{#1}}

\begin{document}
\hypersetup{pageanchor=false}
\begin{titlepage}
\centering
\vspace*{0.35in}
{\large\scshape Research article with complete written proofs\par}
\vspace{0.45in}
{\Huge\bfseries\color{ink} Forcing, Omitted Cuts,\par and Old Surreal Fields\par}
\vspace{0.2in}
{\Large Ordinal-sequence spectra, surcomplex conjugation,\par and a Boolean completion of reverse simplicity\par}
\vspace{0.35in}
{\large Prepared with ChatGPT\par}
{22 September 2026\par}
\vspace{0.3in}
\begin{minipage}{0.94\textwidth}
\begin{abstract}
Let $M\subseteq N$ be transitive universes of ZFC with the same ordinals,
with the ground universe available as a definable predicate, and let
$K=\No^M$ be the old surreal field viewed from $N$. We prove that, for every
uncountable cardinal $\kappa$ of $N$, the ordered field $K$ is
$\kappa$-saturated in $N$ exactly when $N$ has no new ordinal-valued sequences
of length less than $\kappa$ over $M$. In fact, the latter condition is
equivalent to interpolation for every set-sized cut with just one side of
size less than $\kappa$. The converse is established by a set-sized,
arbitrarily branching tree of nested surreal intervals with an explicit
branch decoder. A first new sequence of length $\delta$ produces a gap of
exact cofinality pair $(\delta,\delta)$; $\delta$ is regular in both universes.

The countable boundary is different: $K$ is $\omega$-saturated exactly when
no new reals are added. In contrast, the pure algebraically closed field
$K[i]$ remains saturated over every set of parameters, regardless of new
sequences. Naming conjugation restores precisely the saturation spectrum
of $K$. We also prove a sharp preservation criterion for bounded-birthday
fields, explain Cohen, closed-forcing and Prikry examples, and show why
these new order gaps do not produce nonconstant set-indexed Cauchy nets.
Finally, we explicitly construct a set-complete class Boolean algebra in
which the reverse-simplicity order on surreal signs is dense, answering
a precisely interpreted question in Berenbeim's 2020 notes. Completeness
for all subclasses is ruled out under Global Choice.
\end{abstract}
\end{minipage}
\vfill
\begin{minipage}{0.94\textwidth}\small
\textbf{Status and priority.} The principal spectrum and decoding results
are proposed contributions with proofs, not a claim of certified historical
priority. They were not located in the consulted sources or the specified
repository review. Classical surreal facts, quantifier elimination, and
standard forcing facts are attributed separately. No Lean verification or
independent referee review is claimed. Repository snapshot: \shorthash.
\end{minipage}
\end{titlepage}
\hypersetup{pageanchor=true}
\pagenumbering{roman}
\setcounter{tocdepth}{2}
{\small\setlength{\parskip}{0pt}\tableofcontents}
\clearpage
\pagenumbering{arabic}

\section{The questions and the answers}\label{sec:overview}

Forcing preserves ordinals, but it can add new functions on old ordinals.
Consequently, it can add new surreal sign sequences without changing the
ordinal labels used for birthdays. This observation raises a more delicate
question than whether new surreal numbers exist:
\begin{quote}
How much of the old surreal field's ability to realize cuts and first-order
types survives when the surrounding universe acquires new sets?
\end{quote}
An old field can remain an elementary subfield of a larger surreal field
while losing saturation. Elementarity tests individual first-order formulas
with finitely many parameters. Saturation additionally tests entire sets of
formulas, and those sets themselves can be new.

Throughout the main argument, $M\subseteq N$ have the same ordinals, and
\[
 K=\No^M,\qquad K^+=\No^N,\qquad C=K[i],\qquad C^+=K^+[i].
\]
All cardinalities, types, and set-sized families called \emph{external}
are computed in $N$. They are not computed in a still larger metatheory.
Section~\ref{sec:foundations} makes the class conventions precise.

For an infinite cardinal $\kappa$ of $N$, write
\begin{equation}\label{eq:NSdef}
 \NS_{<\kappa}(M,N)
 \quad\Longleftrightarrow\quad
 ({}^\lambda\theta)^N\subseteq M
 \quad\text{for all ordinals }\lambda<\kappa\text{ and }\theta.
\end{equation}
Every ordinal-valued sequence has bounded range, so this is exactly the
absence of new ordinal-valued sequences of length less than $\kappa$.
The range bound $\theta$ is \emph{not} required to be below $\kappa$.

\begin{theorem}[Main spectrum theorem]\label{thm:main-summary}
For every infinite cardinal $\kappa$ of $N$, the following are equivalent.
\begin{enumerate}[label=\textnormal{(\roman*)}]
\item $\NS_{<\kappa}(M,N)$ holds.
\item Every pair of $N$-sets $L,R\subseteq K$ with $L<R$ and
$\min(\abs L,\abs R)<\kappa$ has a separator in $K$.
\item Every such cut with $\abs{L\cup R}<\kappa$ has a separator in $K$.
\end{enumerate}
For uncountable $\kappa$, these are also equivalent to $\kappa$-saturation
of $K$ in the ordered-field language, as evaluated in $N$.
\end{theorem}

Condition (ii) is stronger in appearance than the usual small-cut condition:
the other side can be an arbitrarily large set. Neither side may be a proper
class. An especially useful special case is that a cut is always filled if
one side is an actual ground-model set.

\begin{theorem}[First loss and language sensitivity]\label{thm:package-summary}
Suppose $M\ne N$. Let $\delta$ be the least ordinal length of a new
ordinal-valued sequence. Then:
\begin{enumerate}[label=\textnormal{(\roman*)}]
\item $\delta$ is an infinite regular cardinal in both $M$ and $N$.
\item The order $K$ has a gap with set-sized cofinal and coinitial witnesses
of exact cofinality pair $(\delta,\delta)$. No such gap has either
cofinality smaller than $\delta$.
\item For uncountable $N$-cardinals $\kappa$, $K$ is $\kappa$-saturated
exactly when $\kappa\leq\delta$.
\item $K$ is $\omega$-saturated exactly when $\R^M=\R^N$.
\item $C$ in the pure field language is saturated over every $N$-set of
parameters. The expansion $(C,c)$ by $c(a+bi)=a-bi$ has exactly the
saturation spectrum of $K$, including at $\omega$.
\end{enumerate}
\end{theorem}

The article also treats the set-sized field
$F_\kappa=(\No_{<\kappa})^M$. For an uncountable regular ground cardinal
$\kappa$ that remains a cardinal in $N$, it proves
\begin{equation}\label{eq:bounded-summary}
\begin{split}
 F_\kappa\text{ is }\kappa\text{-saturated in }N
 \quad\Longleftrightarrow\quad&
 \cf^N(\kappa)=\kappa\ \text{ and}\\[-2pt]
 &({}^\alpha 2)^N=({}^\alpha 2)^M\quad(\alpha<\kappa).
\end{split}
\end{equation}
Thus the full-class theorem and the bounded-fragment theorem do not test
identical preservation properties.

\subsection{What is being offered as new}

The proposed contributions are the exact external spectrum in
Theorem~\ref{thm:main-summary}, its explicit arbitrary-alphabet interval
coding, the exact gap threshold, the distinction between $\omega$- and
$\omega_1$-saturation under no-real extensions, and the resulting
surcomplex language comparison. Their ingredients include standard
cut realization and quantifier elimination; those ingredients are not
claimed as new. Several consequences are deliberately labeled consequences
rather than independent discoveries.

A separate section answers a concrete question posed in Berenbeim's
2020 \emph{notes}, not a conjecture attributed to a peer-reviewed paper.
Its construction is a natural direct-limit construction and may be
folklore in class forcing. Its value here is an explicit, size-correct
answer, not a claim that Boolean completions were previously unavailable.
General class-forcing completion theory predates those notes
\cite{HolyEtAl}.

Berenbeim also explicitly proposes studying surreals under forcing in the
final section of the same notes \cite[pp.~32, 34]{Berenbeim}. The present
article develops one concrete direction of that program. It does not claim
to settle the notes' broader suggestions about embeddings of universes.

\section{Foundations, notation, and imported facts}\label{sec:foundations}

\subsection{Two universes, rather than an unspecified outside observer}

A precise interpretation is to work in $N\models\mathrm{ZFC}$ with a
parameter-definable transitive inner class $M$ containing all ordinals and
satisfying ZFC. The assertion that $M$ satisfies ZFC is an axiom scheme.
All uses of Separation and Replacement with the predicate $M$ are then
instances of the ordinary schemes after substituting its definition.
Alternatively, one may use a two-universe metatheory with these closure
properties stipulated. The forcing case $N=M[G]$ is the intended example;
the proofs themselves use only the stated inner-model properties.

The superscript in $\No^M$ means that the sign functions, their construction,
and the operations are computed in $M$. In $N$, $K$ is a proper class:
it contains an old all-plus sign of every ordinal length. It is not an
additional set-sized model assumed to have a cardinality.

If one instead represents $M$ and $N$ by countable transitive \emph{set}
models in a larger universe, all claims about sets, cardinalities and
saturation must still be interpreted internally to $N$. They are not
claims of saturation with respect to every subset in that larger universe.
This distinction is indispensable.

We do not assume an unrestricted truth predicate for arbitrary class
structures. In the two languages used here, quantifier elimination provides
a uniform reduction of formulas to finite algebraic conditions. It therefore
supplies the required satisfaction relation directly. Types are ordinary
$N$-sets of finite syntactic expressions over set-sized parameter sets.
This avoids the general class-satisfaction obstacle discussed, for example,
in the foundational remarks of van den Dries--Ehrlich \cite{vdDE}.

\subsection{Signs and birthdays}

A surreal number is a sign function $s:\alpha\to\{-,+\}$ for an ordinal
$\alpha$; its birthday is $\bd(s)=\alpha$. The empty function represents
zero. At the first place where two signs differ, numerical comparison uses
\[
 -\ <\ \text{termination}\ <\ +.
\]
The simplicity relation is proper initial segment, not numerical order.
The ordinal $\alpha$ is represented by the all-plus function of length
$\alpha$. Its surreal successor agrees with $\alpha+1$ as an ordinal;
arbitrary ordinal addition and surreal field addition must not be conflated.

For sets $L,R$ of surreals, $L<R$ means $\ell<r$ for every
$\ell\in L$ and $r\in R$. Empty sides are allowed. A separator means a
\emph{strict} separator: $L<x<R$. The classical cut theorem supplies the
unique simplest separator, denoted $\cut LR$. Standard references are
Conway, Gonshor, and van den Dries--Ehrlich \cite{Conway,Gonshor,vdDE}.
We use the birthday bound
\begin{equation}\label{eq:cutbirthday}
 \bd\!\left(\cut LR\right)
 \leq\sup\{\bd(y)+1:y\in L\cup R\}.
\end{equation}
The supremum of the empty set is zero. In particular, every set-sized
separated cut in the full surreal class has a separator.

For clarity, the familiar sign criterion behind several later arguments
is recorded explicitly.
\begin{lemma}[The cone of a sign]\label{lem:signcone}
For $s:\alpha\to\{-,+\}$ put
\[
 L_s=\{s\res\beta:\beta<\alpha,\ s(\beta)=+\},\qquad
 R_s=\{s\res\beta:\beta<\alpha,\ s(\beta)=-\}.
\]
A surreal $x$ satisfies $L_s<x<R_s$ if and only if $s$ is an initial
segment of $x$, allowing $s=x$.
\end{lemma}
\begin{proof}
If $x$ extends $s$, its comparison with each prefix in the displayed sets
is decided by the next sign of $s$, so the required strict inequalities hold.
Conversely, if $x$ terminates at $\beta<\alpha$ while agreeing up to that
point, then $x=s\res\beta$ violates the corresponding strict inequality.
If it first disagrees at $\beta<\alpha$, compare $x$ with the terminating
prefix $s\res\beta$. The opposite signs place $x$ on the wrong side of
that prefix. These exhaust the ways in which $x$ could fail to extend $s$.
\end{proof}

The numerical bounds
\begin{equation}\label{eq:birthbounds}
 -\bd(x)\leq x\leq\bd(x)
\end{equation}
follow immediately by comparing a sign with the all-minus and all-plus
signs of the same length. The ordinals in this formula are their surreal
images.

\subsection{Field theory and types}

The full surreal field is real closed, and its extension by a square root
of $-1$ is algebraically closed. Also, $\No_{<\lambda}$ is a real closed
subfield whenever $\lambda$ is an epsilon number. In particular, this holds
for every uncountable cardinal $\lambda$: all such cardinals satisfy
$\omega^\lambda=\lambda$ in ordinal arithmetic. We use the corrected
birthday-fragment results of van den Dries--Ehrlich and the explicit
restatement in Bournez--Guilmant \cite{vdDE,vdDEerr,BournezGuilmant}.
Nothing about closure under exponential functions is needed below.

We use quantifier elimination for real closed fields in
$\Lang_{\mathrm{or}}=\{0,1,+,-,\cdot,<\}$ and for algebraically closed
fields of characteristic zero in $\Lang_{\mathrm{fld}}=\{0,1,+,-,\cdot\}$
\cite{Marker}. Two concrete consequences are worth stating.
First, an embedding between real closed ordered fields is elementary;
an embedding between algebraically closed fields of the same characteristic
is elementary. Second, over a real closed parameter subfield a nonalgebraic
one-variable type is determined by an order cut, while over an algebraically
closed field one-variable types are algebraic or transcendental.

Our convention is that a structure is $\kappa$-saturated if every complete
one-variable type over a parameter set of cardinality less than $\kappa$
that is finitely satisfiable in the structure is realized there. As usual,
this is equivalent to the finite-variable convention. One proves that
extension by induction on the number of variables, using compactness and
then naming each chosen realization. This induction also works for the
class structures here: all syntactic data and parameter sets at every step
are sets, and satisfaction is supplied by quantifier elimination.

For an order, we use \emph{$\eta_\kappa$} to mean that every separated cut
with $\abs{L\cup R}<\kappa$ has a separator. This cardinal-index notation
is defined here to avoid the alternative convention in which $\eta_\alpha$
is indexed by an aleph subscript.

\begin{lemma}[Real closed saturation from cuts]\label{lem:RCFsat}
For an uncountable cardinal $\kappa$, a real closed ordered field is
$\kappa$-saturated if and only if its order is $\eta_\kappa$.
The assertion applies also to the definable class fields used here.
\end{lemma}
\begin{proof}
Saturation realizes the finitely satisfiable partial type consisting of
the inequalities of a small cut; a complete extension exists by ordinary
compactness applied to the set of formulas with those parameters.

Conversely, let $A$ have size less than $\kappa$. The real closure $F$ of
$\Q(A)$ inside the field is a set of cardinality at most
$\max(\abs A,\aleph_0)<\kappa$. A complete type over $A$ can be extended
to one over $F$ which is finitely satisfiable in the ambient field: append
the elementary diagram on $F$ and apply compactness. If the extended type
is algebraic, it is realized in $F$, since $F$ is real closed. Otherwise
it declares, for every $a\in F$, exactly one of $a<X$ and $X<a$. These
inequalities form a separated cut of size less than $\kappa$.
A separator realizes the entire type by quantifier elimination, because
all polynomial signs over $F$ are determined by its position relative to
the real roots of those polynomials. Restrict back to $A$.
\end{proof}

The uncountability hypothesis cannot be dropped. Every dense order without
endpoints has the finite-cut property $\eta_\omega$, whereas a real closed
field need not be $\omega$-saturated.

\section{Absoluteness and ground-set bounds}\label{sec:bounds}

\begin{proposition}[Old arithmetic and elementary inclusions]\label{prop:absolute}
The sign inclusion $K\subseteq K^+$ preserves birthdays, numerical order,
addition, multiplication, and inversion. In $N$, $K$ is real closed and
$C$ is algebraically closed. Moreover,
\[
 K\preccurlyeq K^+\quad\text{in }\Lang_{\mathrm{or}},\qquad
 C\preccurlyeq C^+\quad\text{in }\Lang_{\mathrm{fld}}.
\]
\end{proposition}
\begin{proof}
Old sign functions have the same domains and entries, so birthdays and
first-disagreement comparisons are absolute. Their canonical option sets
consist of their proper prefixes and are identical in the two universes.
For an old cut presentation, its simplest separator is likewise absolute.
Indeed, its old simplest separator $s$ exists. In $N$, let $t$ be the
newly computed simplest separator. The simplicity property of cuts says
that $t$ is an initial segment of every separator, in particular of $s$.
Thus $t$ is old. Since the old construction already chose the simplest
old separator, $t=s$. Alternatively, the standard
sign-by-sign construction of the simplest separator is a set recursion
with the same old input and uniquely determined successive signs.

Conway's addition and multiplication are defined from canonical options
by well-founded, set-like recursion. Induction on the recursive arguments
therefore gives the same values in the two universes. Every recursion
needed for particular inputs has a set-sized dependency closure; no new
class-recursion principle is being assumed. The inverse of an old nonzero
number remains its inverse by the already absolute multiplication law
and uniqueness.

A polynomial over $K$ has finitely many old coefficients. That finite
coefficient tuple, and hence the polynomial, belongs to $M$. In $M$,
positive elements have square roots and odd-degree polynomials have roots.
Those same witnesses work in $N$. Thus $K$ is externally real closed.
The corresponding finite-polynomial argument proves external algebraic
closedness of $C$. Quantifier elimination now gives both elementary
inclusions.
\end{proof}

This proposition asserts elementarity of the \emph{fields}. It does not
assert $M\preccurlyeq N$ as models of set theory, which is generally false
for a forcing extension.

\begin{lemma}[Every external set of old numbers is ground-bounded]\label{lem:groundbound}
If $A$ is an $N$-set and $A\subseteq K$, there are an ordinal $\alpha$
and a ground set $B\in M$ with
\[
 A\subseteq B=(\No_{<\alpha})^M.
\]
Every $N$-set of elements of $C$ is likewise contained in a ground set.
\end{lemma}
\begin{proof}
In $N$ form $\alpha=\sup\{\bd(a)+1:a\in A\}$. This ordinal belongs to
$M$, because the universes have the same ordinals. The set of all old signs
of lengths below $\alpha$ exists in $M$ by Power Set and Replacement,
and it contains every $a\in A$. For $C$, do the same for both coordinates
and take the ground set of their ordered pairs.
\end{proof}

The lemma is not a covering property of a forcing notion. It holds even
when a forcing badly fails familiar small-set covering properties. The
containing ground set can be vastly larger than $A$.

\begin{lemma}[Coding short families]\label{lem:shortold}
Assume $\NS_{<\kappa}(M,N)$. Every $N$-set of old surreal or surcomplex
numbers of cardinality less than $\kappa$ is a member of $M$. More generally,
every sequence of length less than $\kappa$ with values in a fixed ground
set is a ground sequence.
\end{lemma}
\begin{proof}
For a ground set $B$, fix in $M$ an injection $e:B\to\theta$ for an
ordinal $\theta$, using Choice in $M$. Composing a short sequence in $B$
with $e$ gives an ordinal-valued sequence of the same length, hence an old
one. Decoding with $e$ shows that the original sequence is old. For an
external set $A$ of size less than $\kappa$, use Lemma~\ref{lem:groundbound}
and choose in $N$ an enumeration of $A$ of length less than $\kappa$.
The enumeration is old, and its range belongs to $M$.
\end{proof}

The following asymmetric fact is the efficient direction of the main theorem.
\begin{theorem}[One old side suffices]\label{thm:oneside}
Let $L,R$ be $N$-sets contained in $K$, with $L<R$. If either $L\in M$
or $R\in M$, then some $x\in K$ satisfies $L<x<R$.
\end{theorem}
\begin{proof}
Suppose $L\in M$. Choose a ground set $B\supseteq R$ by
Lemma~\ref{lem:groundbound}, and define, in $M$,
\[
 U=\{b\in B:(\forall\ell\in L)\ \ell<b\}.
\]
Order is absolute, so $R\subseteq U$ and $L<U$. The ground cut
$x=\cut L U$, constructed in $M$, satisfies $L<x<U$, and therefore
$L<x<R$. This also covers $R=\varnothing$ and $L=\varnothing$.
If $R\in M$, apply the proved case to $-R<-L$ and negate its separator.
\end{proof}

\begin{corollary}\label{cor:NSoneside}
If $\NS_{<\kappa}(M,N)$ holds, every external set-sized cut with at least
one side of size less than $\kappa$ is filled in $K$.
\end{corollary}
\begin{proof}
Its small side belongs to $M$ by Lemma~\ref{lem:shortold}. Apply
Theorem~\ref{thm:oneside}.
\end{proof}

In particular, every $N$-set of positive elements of $K$ has a positive
lower bound in $K$: take $L=\{0\}$. Every $N$-set of old numbers also
has an old strict upper bound, either by the theorem with $R=\varnothing$
and the symmetric case, or directly by bounding birthdays and using a
larger all-plus ordinal. These statements survive every extension under
consideration, including ones that destroy countable saturation.

\section{Coding a new sequence by nested old intervals}\label{sec:decoder}

The converse requires more than the observation that a new binary sign
has a new cut. A short sequence can take values in a very large ordinal.
For example, a new countable cofinal sequence can appear while no real is
added. Encoding it as one binary sign might require a long ordinal domain,
which would destroy the desired cardinal bound on the cut.

The construction below keeps the number of inequalities equal to the
length of the sequence, independently of the size of its alphabet.

\begin{lemma}[Arbitrarily branching interval tree]\label{lem:tree}
Let $\theta\geq2$ be an ordinal and let $\lambda$ be a nonzero limit
ordinal, both in $M$. In $M$ there is a set-sized table of pairs
\[
 (a_s,b_s)\in K^2,
 \qquad s\in T:=\bigl({}^{<\lambda}\theta\bigr)^M,
\]
with these properties:
\begin{enumerate}[label=\textnormal{(\roman*)}]
\item $a_\varnothing=0<b_\varnothing=1$.
\item If $s$ is a proper initial segment of $t$, then
$a_s<a_t<b_t<b_s$.
\item The intervals belonging to distinct immediate successors of one
node are pairwise disjoint, and are ordered by their successor labels.
\end{enumerate}
\end{lemma}
\begin{proof}
Perform recursion on the set of levels below $\lambda$. Suppose an interval
$(a_s,b_s)$ is given and write $w_s=b_s-a_s>0$. Regard $\xi$ and $\theta$
as their all-plus surreal images and put
\begin{equation}\label{eq:split}
\begin{split}
 a_{s^\frown\xi}&=a_s+w_s\frac{3\xi+1}{3\theta+3},\\
 b_{s^\frown\xi}&=a_s+w_s\frac{3\xi+2}{3\theta+3}
 \qquad(\xi<\theta).
\end{split}
\end{equation}
Every operation in this display is \emph{surreal field arithmetic}.
If $\xi<\eta<\theta$, the equality of ordinal and surreal successor
implies $\xi+1\leq\eta$. Hence
\[
 0<3\xi+1<3\xi+2<3\eta+1<3\eta+2<3\theta+3.
\]
This proves the required interior and sibling-separation properties.

At a nonzero limit level $\beta<\lambda$, let $s\in({}^\beta\theta)^M$.
The induction hypothesis gives separated ground sets
\[
 A_s=\{a_{s\res\gamma}:\gamma<\beta\},\qquad
 B_s=\{b_{s\res\gamma}:\gamma<\beta\}.
\]
Define successively, inside $M$,
\begin{equation}\label{eq:limitinterval}
 a_s=\cut{A_s}{B_s},\qquad
 b_s=\cut{A_s\cup\{a_s\}}{B_s}.
\end{equation}
Then $A_s<a_s<b_s<B_s$. This verifies strict nesting even at limit levels.
Each level and its prefix data are sets in $M$; Replacement and ordinary
set-length recursion produce the entire table in $M$. No recursion along
all ordinals is used.
\end{proof}

\begin{theorem}[Branch decoder and omitted cut]\label{thm:decode}
Suppose $f:\lambda\to\theta$ belongs to $N\setminus M$, where $\lambda$
is a nonzero limit ordinal, and every proper initial segment $f\res\alpha$
belongs to $M$. Then $K$ omits the finitely satisfiable partial type
\begin{equation}\label{eq:branchtype}
 p_f(X)=\{\ a_{f\res\alpha}<X<b_{f\res\alpha}:\alpha<\lambda\ \}.
\end{equation}
It uses at most $\abs\lambda$ parameters on each side. The left endpoints
strictly increase, and the right endpoints strictly decrease.
\end{theorem}
\begin{proof}
All indicated endpoints are old. Any finite subset of the displayed
inequalities is satisfied by a point in its deepest interval, for example
its midpoint, which is an old field element. Strict nesting proves both
monotonicity assertions.

Suppose, toward a contradiction, that an old $x$ realizes all inequalities.
Using only the ground interval table and the old parameter $x$, define
$d_x:\lambda\to\theta$ by recursion in $M$. At a successor step, having
constructed a prefix $s$, choose the unique $\xi<\theta$ for which
\[
 a_{s^\frown\xi}<x<b_{s^\frown\xi},
\]
if such a label exists; otherwise choose $0$. At limit steps take unions
of the previously constructed prefixes. Distinct sibling intervals are
disjoint, so the successful choice is unique. The default value makes
this a \emph{total ground recursion}, without presupposing that a branch
through the intervals belongs to $M$.

We prove in $N$ by induction that $d_x\res\alpha=f\res\alpha$ for every
$\alpha<\lambda$. At the next step, $x$ lies in the interval of the true
successor $f\res(\alpha+1)$ by the type. The unique chosen label must
therefore be $f(\alpha)$. At limits, agreement follows by union.
Here $\alpha+1<\lambda$ for each $\alpha<\lambda$, since $\lambda$ is
limit, so every required successor interval is in the table. Thus
$d_x=f$, contradicting $d_x\in M$ and $f\notin M$.
\end{proof}

The proof does not infer that an external branch is old merely because
all its nodes are old. That inference would be false. Instead, an assumed
\emph{old realizer} supplies a total, deterministic ground decoder for the
entire branch. This is the step that makes the argument work.

There is also no assertion that the interval widths tend to zero in the
fine surreal topology. Section~\ref{sec:topology} proves that these widths
actually have a common positive old lower bound.

\section{The exact spectrum and its first gap}\label{sec:spectrum}

\subsection{Proof of the main equivalence}

\begin{proof}[Proof of Theorem~\ref{thm:main-summary}]
The implication (i)$\Rightarrow$(ii) is Corollary~\ref{cor:NSoneside};
(ii)$\Rightarrow$(iii) is immediate.

Suppose (i) fails. Choose the least ordinal $\lambda<\kappa$ for which
there is a new ordinal-valued sequence of length $\lambda$. Choose one
such sequence and bound its range by an ordinal $\theta\geq2$. No finite
sequence of old ordinals is new. Also $\lambda$ cannot be a successor:
a sequence at a successor stage consists of its shorter prefix and one
old ordinal, and both would already belong to $M$. Hence $\lambda$ is
an infinite limit ordinal, and all proper prefixes of the chosen sequence
are old. Theorem~\ref{thm:decode} produces an unfilled cut whose two sides
together have cardinality at most $\abs\lambda<\kappa$ in $N$. Thus
(iii) fails. This proves (iii)$\Rightarrow$(i).

For uncountable $\kappa$, Proposition~\ref{prop:absolute} and
Lemma~\ref{lem:RCFsat} identify (iii) with the asserted saturation property.
\end{proof}

Notice that $\kappa$ need not be regular. The theorem is stated for every
infinite $N$-cardinal; only the conversion to real closed saturation
requires $\kappa>\omega$.

\subsection{The least new sequence has regular length}

\begin{definition}
When $M\ne N$, define the \emph{first sequence length}
\begin{equation}\label{eq:delta}
 \delta(M,N)=\min\{\lambda:\exists\theta\ 
 ({}^\lambda\theta)^N\not\subseteq M\}.
\end{equation}
\end{definition}

This minimum exists. Indeed, if there were no new sets of ordinals, every
set in $N$ could be coded there by a well-founded extensional relation on
an ordinal, together with a distinguished point. Such a relation is itself
coded by a set of ordinals. If the code is in $M$, its Mostowski collapse
in $M$ agrees with the collapse in $N$, yielding the original set in $M$.
Thus absence of all new sets of ordinals would imply $M=N$. A new set of
ordinals gives a new characteristic sequence.

\begin{proposition}\label{prop:deltaregular}
The ordinal $\delta(M,N)$ is an infinite regular cardinal in both $M$ and $N$.
\end{proposition}
\begin{proof}
Write $\delta=\delta(M,N)$ and choose a new $f:\delta\to\theta$.
We have already observed that $\delta$ is infinite and all its proper
prefixes are old.

If $\delta$ were not an $N$-cardinal, choose in $N$ a bijection
$e:\mu\to\delta$ with $\mu<\delta$. By minimality of $\delta$, $e\in M$.
The sequence $f\circ e$ also has length less than $\delta$, hence belongs
to $M$. Decoding with $e$ would give $f\in M$, a contradiction. Thus
$\delta$ is a cardinal in $N$, and therefore also in $M$.

If $\cf^N(\delta)=\mu<\delta$, choose a cofinal sequence
$e:\mu\to\delta$. It belongs to $M$ by minimality. Every
$f\res e(\xi)$ is old and lies in the fixed ground set
\[
 T=({}^{<\delta}\theta)^M.
\]
Well-order $T$ in $M$. The sequence of the corresponding ordinal codes
of $f\res e(\xi)$, for $\xi<\mu$, has length less than $\delta$ and
therefore is old. Its ground decoding and union reconstruct $f$, again a
contradiction. Hence $\delta$ is regular in $N$. Any shorter ground
cofinal sequence would still be cofinal in $N$, so $\delta$ is regular
in $M$ as well.
\end{proof}

This is an elementary minimal-length fact about extensions, not a claim
of a new characterization of regular cardinals in general.

\subsection{Exact two-sided cofinalities}

A \emph{gap} in the class order $K$ is a partition $(D,E)$ into a lower
class and an upper class with $D<E$, neither having an endpoint toward the
other. We call it a \emph{set-character gap} if $D$ has a set-sized
cofinal subset and $E$ has a set-sized coinitial subset. Its cofinality
pair is the least possible pair of cardinalities of such witnesses,
computed in $N$. We never require $D$ or $E$ themselves to be sets.

\begin{theorem}[First gap theorem]\label{thm:firstgap}
Let $M\ne N$ and $\delta=\delta(M,N)$. Every set-character gap of $K$
has both cofinalities at least $\delta$, and $K$ has one of exact
cofinality pair $(\delta,\delta)$.
\end{theorem}
\begin{proof}
Apply Theorem~\ref{thm:decode} to a new sequence of minimal length
$\delta$, and abbreviate its endpoints as $a_\alpha,b_\alpha$.
Define in $N$
\[
\begin{split}
 D&=\{x\in K:\exists\alpha<\delta\ (x\leq a_\alpha)\},\\
 E&=\{x\in K:\exists\alpha<\delta\ (b_\alpha\leq x)\}.
\end{split}
\]
All left endpoints are below all right endpoints. A point in neither
class would satisfy every strict interval inequality, which is impossible.
Thus $D\cup E=K$ and $D<E$. Strict monotonicity and the fact that
$\delta$ is limit show that neither class has an endpoint toward the
other. The displayed endpoint sequences are cofinal and coinitial
witnesses of size $\delta$.

For the lower bound, consider any set-character gap with witnesses $A$
on the lower side and $B$ on the upper side. If either witness could have
size less than $\delta$, Theorem~\ref{thm:main-summary}, applied at
$\kappa=\delta$, would give an old strict separator between $A$ and $B$.
A point of the lower class cannot lie strictly above a cofinal subset
of that class, and a point of the upper class cannot lie strictly below
a coinitial subset of that class. Since the two classes partition $K$,
such a separator is impossible. Both cofinalities are therefore at least
$\delta$, establishing equality for the constructed gap.
\end{proof}

\begin{corollary}[Recovering the sequence threshold from the order]\label{cor:recover}
For $M\ne N$, the ordinal $\delta(M,N)$ is the least cardinal $\mu$ for
which $K$ has an unfilled set-sized cut with both sides of size at most
$\mu$. Equivalently, it is the least cofinality appearing on either side
of a set-character gap of $K$.
\end{corollary}
\begin{proof}
The upper bounds follow from Theorem~\ref{thm:firstgap}. A smaller
unfilled cut contradicts the main theorem at $\delta$, and a smaller
gap cofinality contradicts the first-gap theorem.
\end{proof}

\begin{corollary}\label{cor:spectrum}
For uncountable $N$-cardinals $\kappa$, the old ordered field $K$ is
$\kappa$-saturated exactly when $\kappa\leq\delta(M,N)$. If
$\delta=\omega$, no uncountable level of saturation survives.
\end{corollary}
\begin{proof}
There are no new sequences of length less than $\delta$, and there is
one of length $\delta$. Apply the main theorem.
\end{proof}

The statement that $K$ has no smaller \emph{set-character} gap does not
exclude cuts with proper-class cofinality on one side. Those are outside
the invariant defined here and outside the main interpolation theorem.

\section{The countable boundary: no reals is weaker}\label{sec:omega}

The absence of new reals and the absence of new countable ordinal-valued
sequences are different preservation properties. The old surreal field
records this distinction exactly.

\begin{theorem}[Finite-parameter saturation]\label{thm:omega}
The following are equivalent:
\[
 K\text{ is }\omega\text{-saturated in }N
 \quad\Longleftrightarrow\quad
 \Pow(\omega)^M=\Pow(\omega)^N
 \quad\Longleftrightarrow\quad
 \R^M=\R^N.
\]
\end{theorem}
\begin{proof}
Equality of the two power sets is equivalent to equality of the reals by
ordinary countable coding. We prove the model-theoretic equivalence.

Suppose no new subsets of $\omega$ are added. Let $A$ be a finite set
of old parameters, and let $p$ be a complete one-variable type over $A$
which is finitely satisfiable in $K$ according to $N$. The finite set
$A$ belongs to $M$. The set of formulas with parameters in $A$ has a
countable enumeration in $M$. The code of $p$ under this enumeration
is a subset of $\omega$, so $p\in M$.

For each finite subset $p_0$ of $p$, the assertion that its conjunction
has a solution is an existential ordered-field formula over finitely
many old parameters. Proposition~\ref{prop:absolute} and quantifier
elimination show that this assertion has the same truth value in the
ground and external interpretations of $K$. Hence $M$ recognizes $p$
as finitely satisfiable. In $M$, the full surreal field realizes every
set-sized type over a set of parameters: take a set real closure of the
parameter field, use the cut theorem, and apply the argument of
Lemma~\ref{lem:RCFsat}. Thus $p$ has an old realization.

Conversely, let $r\in\R^N\setminus\R^M$. Consider the rational-cut
partial type
\begin{equation}\label{eq:newrealtype}
 \{q<X:q\in\Q,\ q<r\}
 \ \cup\
 \{X<q:q\in\Q,\ r<q\}.
\end{equation}
It is finitely satisfiable by rational numbers. Rational constants can
be expressed without parameters in the field language: for example,
$m/n<X$ is equivalent to $m<nX$ for $n>0$. Thus this is a partial type
over the empty parameter set.

If an old $x$ realized it, the ground set
$\{q\in\Q:q<x\}$, formed by Separation in $M$, would be exactly the
lower rational cut of $r$. Ground Dedekind completeness would reconstruct
$r$ in $M$, a contradiction. For a complete omitted type, take the type of the canonical surreal
image of $r$ in $K^+$ over the empty parameter set. This is an $N$-set
of formulas, obtained using quantifier elimination, and it contains
\eqref{eq:newrealtype}. Every finite part is realized in $K$ because
$K\preccurlyeq K^+$ by Proposition~\ref{prop:absolute}. Thus $K$ is
not $\omega$-saturated. No class compactness principle is required.
\end{proof}

\begin{corollary}\label{cor:omegaone}
An extension which adds no reals but adds a countable ordinal-valued
sequence leaves $K$ $\omega$-saturated and destroys
$\omega_1^N$-saturation. It creates a $(\omega,\omega)$ set-character
gap in $K$, but none of its omitted types can be witnessed over a fixed
finite parameter set.
\end{corollary}
\begin{proof}
The first new sequence length is $\omega$. Apply
Theorems~\ref{thm:omega} and \ref{thm:firstgap} and
Corollary~\ref{cor:spectrum}. The last assertion is precisely the
surviving $\omega$-saturation.
\end{proof}

The new countable parameter family in the interval-tree construction
carries information that cannot be coded by an old finite parameter tuple.
This is why a bare new-real argument would miss this example.

\section{Surcomplex numbers: a language-dependent dichotomy}\label{sec:complex}

Write $C=K[i]$ and let $c(a+bi)=a-bi$. The repository's automorphism
report already emphasizes the distinction between a pure algebraically
closed field and a field with a distinguished surreal real axis
\cite{RepoAutomorphisms}. That distinction is standard algebra and is
not claimed as new here. The point of this section is its exact effect
on \emph{external saturation under a change of universe}.

\subsection{The pure field forgets the sequence obstruction}

\begin{theorem}[Absolute set-saturation of the old pure surcomplex field]
\label{thm:purecomplex}
For every $N$-set $A\subseteq C$, every complete one-variable pure-field
type over $A$ which is finitely satisfiable in $C$ is realized in $C$.
Consequently, $C$ is $\kappa$-saturated in $N$ for every $N$-cardinal
$\kappa$, with no assumption on new sequences or new reals.
\end{theorem}
\begin{proof}
Choose a ground set $B\supseteq A$ by Lemma~\ref{lem:groundbound}.
In $M$, let
\[
 E=\{z\in C:z\text{ is algebraic over }\Q(B)\}.
\]
This is a ground set. The polynomials over the set field $\Q(B)$ form
a set, each nonzero polynomial has finitely many roots, and Replacement
followed by Union collects those roots. Since $C$ is algebraically
closed, $E$ is an algebraically closed ground subfield.

Choose an old $t\in C\setminus E$. Such a choice does not require a
class-sized transcendence basis. For instance, bound the birthdays of
the real and imaginary coordinates of all elements of the ground set
$E$, and take a larger all-plus ordinal. Then $t$ is transcendental over
$\Q(B)$. This remains true in $N$, since a polynomial witnessing
algebraicity would have a finite old coefficient tuple and would
already be a ground polynomial. In particular, $t$ is transcendental
over $\Q(A)$ as formed in $N$, because every finite rational expression
in $A\subseteq B$ lies in the old field $\Q(B)$.

Now let $p$ be the given complete type. If $p$ contains no equality
$P(X)=0$ for a nonzero polynomial $P$ over $\Q(A)$, completeness and
quantifier elimination imply that $p$ is the unique transcendental
type over $\Q(A)$. It is realized by $t$.

Otherwise choose a nonzero polynomial equation $P(X)=0$ in $p$.
Its finitely many roots $r_1,\ldots,r_n$ all lie in $C$ by
Proposition~\ref{prop:absolute}. If none realized $p$, choose for each
$r_j$ a formula $\varphi_j(X)\in p$ that $r_j$ fails. Then
\[
 \{P(X)=0,\varphi_1(X),\ldots,\varphi_n(X)\}
\]
is a finite subset of $p$ with no solution in $C$, a contradiction.
Thus one of the roots realizes $p$. This argument also covers algebraic
parameters and any new organization of the parameter set $A$.
\end{proof}

There is no cardinal paradox. The field $C$ is a proper class in $N$;
the statement is a scheme over set-sized parameter sets, not the claim
that a set of cardinality $\mu$ realizes types over itself of size $\mu$.
Nor does the theorem by itself produce a class isomorphism between $C$
and $C^+$. A global isomorphism would require a separate class back-and-forth
argument and appropriately stated class principles.

\subsection{Naming conjugation restores the old ordered field}

\begin{lemma}[Real coordinates are interpretable]\label{lem:coordinates}
The fixed field of $c$ is definable in $(C,c)$ by $c(x)=x$. On that
field, the order is definable by
\begin{equation}\label{eq:orderdef}
 x<y\quad\Longleftrightarrow\quad
 \exists u\,\bigl(c(u)=u\ \wedge\ u\ne0\ \wedge\ u^2=y-x\bigr).
\end{equation}
After naming $i$, the map $z\mapsto(\operatorname{re}z,\operatorname{im}z)$
is a definable bijection from $C$ to $K^2$, where
\begin{equation}\label{eq:coordinates}
 \operatorname{re}z=\frac{z+c(z)}2,\qquad
 \operatorname{im}z=\frac{z-c(z)}{2i}.
\end{equation}
All field and conjugation operations translate into polynomial operations
on these two coordinates.
\end{lemma}
\begin{proof}
The fixed-point equation gives exactly the zero imaginary coordinate.
In a real closed field the nonzero squares are exactly the positive
numbers, proving \eqref{eq:orderdef}. The coordinate formulas are
immediate from $i^2=-1$, and their inverse is $(a,b)\mapsto a+bi$.
For example,
\[
 (a,b)(u,v)=(au-bv,av+bu),\qquad c(a,b)=(a,-b).
\]
Consequently quantifiers over $C$ can be replaced by pairs of quantifiers
over $K$, and quantifiers over $K$ by quantifiers over the definable fixed
field.
\end{proof}

\begin{theorem}[Conjugation spectrum]\label{thm:conjugation}
For every infinite $N$-cardinal $\kappa$,
\[
 (C,c)\text{ is }\kappa\text{-saturated in }N
 \quad\Longleftrightarrow\quad
 K\text{ is }\kappa\text{-saturated in }N.
\]
For uncountable $\kappa$, this is equivalent to
$\NS_{<\kappa}(M,N)$. For $\kappa=\omega$, it is equivalent to the
absence of new reals.
\end{theorem}
\begin{proof}
Naming a fixed finite tuple of constants does not alter
$\kappa$-saturation for any infinite $\kappa$: a parameter set of size
less than $\kappa$ stays that small after adjoining finitely many
constants. Conversely, a type in the reduct can be extended to a type
over the named constants by compactness, with finite satisfiability,
and then restricted after realization. We may therefore name $i$.

If $K$ is $\kappa$-saturated, translate a type in $(C,c,i)$ by
Lemma~\ref{lem:coordinates}. A parameter set of size less than $\kappa$
gives a set of its real and imaginary coordinates still of size less
than $\kappa$. The translated type uses two field variables, so the
finite-variable form of saturation of $K$ realizes it. The inverse
coordinate map gives the desired realization in $C$.

Conversely, translate a type over a small subset of $K$ into $(C,c)$
by replacing order with \eqref{eq:orderdef}, restricting quantifiers
to the fixed field, and adding $c(X)=X$. Its finite satisfiability is
preserved. Saturation in $(C,c)$ yields a fixed-field realization, hence
a realization in $K$. The last two assertions follow from the previously
proved spectrum theorems.
\end{proof}

In particular, the interval type \eqref{eq:branchtype} remains omitted
in the conjugation expansion when written using the definable fixed
field and definable positive squares. In the pure field language, those
order inequalities are unavailable; its one-variable transcendental
type is realized by Theorem~\ref{thm:purecomplex}.

\subsection{A useful nonisomorphism consequence}

\begin{corollary}\label{cor:noniso}
Assume $M\ne N$. There is no $N$-class order isomorphism from $K$ onto
$K^+$ whose restrictions to sets have set-sized images. Consequently
there is no such field isomorphism, and no such isomorphism
$(C,c)\cong(C^+,c)$.
\end{corollary}
\begin{proof}
The first-gap theorem supplies two $N$-sets of old endpoints with no
separator in $K$. An isomorphism as stated would send them to separated
$N$-sets in $K^+$, whose cut has a separator. Its inverse would give a
separator in $K$, a contradiction. An isomorphism of real closed fields
preserves order because positivity is definable by nonzero squares.
An isomorphism preserving conjugation restricts to an isomorphism of
the definable real fixed fields.
\end{proof}

The set-image condition is automatic for the ordinary definable class
maps or class maps of GBC with Class Replacement. It is stated explicitly
to avoid quantifying over unrestricted external graphs without an
ambient class theory.

\section{A sharp theorem for bounded birthday fields}\label{sec:fragments}

The previous theorems concern all ground birthdays at once. A set-sized
fragment has a new obstruction: its birthdays can become cofinally
approachable by a short sequence. This can happen even if every individual
sign of length below its birthday bound remains old.

Let $\kappa$ be an uncountable regular cardinal in $M$, and assume it
remains a cardinal in $N$. Put
\[
 F_\kappa=(\No_{<\kappa})^M.
\]
It is a real closed set field in $M$, by the classical birthday-fragment
theorem, and remains real closed in $N$ by the finite-coefficient argument.
Regularity in $N$ is \emph{not} assumed in advance.

\begin{theorem}[Bounded-birthday preservation]\label{thm:fragments}
The following conditions are equivalent.
\begin{enumerate}[label=\textnormal{(\roman*)}]
\item $F_\kappa$ is $\kappa$-saturated in $N$ as an ordered field.
\item The order $F_\kappa$ has the $\eta_\kappa$ property in $N$.
\item $\cf^N(\kappa)=\kappa$ and
$({}^\alpha2)^N=({}^\alpha2)^M$ for every $\alpha<\kappa$.
\item $\kappa$ is regular in $N$ and
$F_\kappa=(\No_{<\kappa})^N$ as sets of sign sequences.
\end{enumerate}
\end{theorem}
\begin{proof}
The equivalence (i)$\Leftrightarrow$(ii) is
Lemma~\ref{lem:RCFsat}. The equivalence (iii)$\Leftrightarrow$(iv) is
exactly the representation of surreals by binary sign sequences.

Assume (iv). Let $L,R\subseteq F_\kappa$ be $N$-sets with
$L<R$ and $\abs{L\cup R}<\kappa$. Since $\kappa$ is regular in $N$,
\[
 \beta=\sup\{\bd(y)+1:y\in L\cup R\}<\kappa.
\]
Construct the cut in $N$. By \eqref{eq:cutbirthday} its birthday is
at most $\beta$, so its separator belongs to $(\No_{<\kappa})^N$,
which equals $F_\kappa$. This proves (ii).

For necessity, first suppose $\cf^N(\kappa)=\mu<\kappa$. A cofinal
sequence of ordinals below $\kappa$, viewed as all-plus surreals, is
cofinal in $F_\kappa$ by \eqref{eq:birthbounds}. It has size less than
$\kappa$, yet has no strict upper bound in $F_\kappa$. Together with
an empty right side it witnesses failure of (ii).

Next suppose some new binary function of length less than $\kappa$
exists. Choose one of least such length, say
$s:\alpha\to\{-,+\}$. Its proper prefixes are all old, and they all
have birthday less than $\kappa$. The canonical prefix sets $L_s,R_s$
of Lemma~\ref{lem:signcone} therefore lie in $F_\kappa$ and together
have size at most $\abs\alpha<\kappa$. Any separator, even in the
full new surreal field, must extend $s$. An old separator cannot do so,
since restricting its old sign function to the old ordinal $\alpha$
would put $s$ in $M$. Thus this small cut is unfilled in $F_\kappa$,
again contradicting (ii). This proves (iii).
\end{proof}

\begin{remark}[Why the two criteria differ]
In the full-class theorem, $\NS_{<\kappa}$ bans short sequences with
\emph{arbitrarily large ordinal range}. The fragment theorem bans new
binary signs below a fixed birthday bound and separately preserves the
cofinality of that bound. A short sequence with unbounded range at a much
larger cardinal may be invisible to one small birthday fragment but
visible to the full old field. Conversely, singularizing the bound itself
can destroy fragment saturation without adding any new sign below it.
\end{remark}

\begin{proposition}[Pure complexification of a set fragment]\label{prop:setACF}
Let $\mu=\abs{F_\kappa}^N$. If $\mu$ is uncountable, the algebraically
closed field $F_\kappa[i]$ is $\mu$-saturated in the pure field language
in $N$, regardless of whether the equivalent conditions of
Theorem~\ref{thm:fragments} hold.
\end{proposition}
\begin{proof}
An algebraically closed field of characteristic zero of uncountable
cardinality $\mu$ has transcendence degree $\mu$ over $\Q$: if a
transcendence basis had cardinality $\nu$, its algebraic closure would
have cardinality $\max(\nu,\aleph_0)$. Over a parameter set of size
less than $\mu$, there is still a transcendental element, and all algebraic
types are realized by algebraic closedness. Quantifier elimination gives
$\mu$-saturation. This is the usual dimension argument for algebraically
closed fields, included as a consequence rather than as a new field theorem.
\end{proof}

\section{Forcing examples and distributivity}\label{sec:examples}

The main theorem applies to all the inner-model pairs in its hypotheses.
For set forcing, it translates the usual preservation of short ordinal
sequences into a field-theoretic statement.

\subsection{The forcing-theoretic formulation}

Recall the convention that $\mathbb P$ is \emph{$<\kappa$-distributive}
if the intersection of any family of fewer than $\kappa$ dense open
subsets is dense, or equivalently that its complete Boolean algebra is
$<\kappa$-distributive. With the usual separative interpretation, this is
equivalent to adding no new ordinal-valued sequences of length less than
$\kappa$ \cite{Jech}. Here ``adds no'' refers to every generic extension,
not merely to a chosen generic filter.

The forward direction can be seen by taking, for every coordinate of a
name for a short sequence, the dense set of conditions deciding that
coordinate. A condition in their intersection decides the entire sequence
as a ground sequence. Conversely, in the complete Boolean algebra, choose
maximal antichains refining the proposed dense open sets. A generic choice
of one member of each antichain is a short ordinal-coded sequence. If
all such sequences are ground sequences, Boolean values deciding an
entire choice provide a dense set below all the chosen antichains,
yielding distributivity. This is the standard Boolean form of the
preservation equivalence, not a new forcing theorem.

Thus, for an uncountable ground cardinal $\kappa$, $<\kappa$-distributive
forcing preserves the $\kappa$-saturation of the full old surreal field.
More exactly, whenever $\kappa$ is a cardinal in the extension, the
criterion for that \emph{particular pair of universes} is
$\NS_{<\kappa}(M,M[G])$. A failure of distributivity of a forcing notion
should not be silently replaced by a claim about every one of its generic
extensions; the theorem is formulated pair by pair to avoid this issue.

\subsection{A Cohen real destroys even finite-parameter saturation}

Let $N=M[g]$ be an extension adding a Cohen real. It adds a subset of
$\omega$, so $\delta(M,N)=\omega$. The old field has a
$(\omega,\omega)$ gap and is not even $\omega$-saturated. One can witness
this directly by the rational-cut type \eqref{eq:newrealtype}, without
using the general interval tree.

Nevertheless, $K\preccurlyeq K^+$ remains true, and the pure field
$C$ remains saturated over every $N$-set of parameters. Therefore neither
loss of real closedness nor failure of elementarity explains the omitted
type: the change lies in which countable sets of formulas are available.
The elementary Cohen forcing facts used here are standard
\cite{Jech,Berenbeim}.

\subsection{Closed forcing can place the threshold at any regular cardinal}

Fix an infinite regular cardinal $\rho$ in $M$. Let
\[
 \mathbb P=\operatorname{Fn}_{<\rho}(\rho,2)
\]
be the set of partial functions from $\rho$ to $2$ with domain of size
less than $\rho$, ordered by reverse inclusion. For $\rho=\omega$,
this is the usual finite-condition Cohen forcing. For uncountable
$\rho$, it is $<\rho$-closed: the union of a descending chain of length
less than $\rho$ is a condition, by regularity of $\rho$.

For completeness, this closure prevents new ordinal-valued sequences
of every length $\lambda<\rho$. Given a name, recursively strengthen a
condition to decide one coordinate at a time; at limit stages and at the
end, closure provides a common extension. Thus the sequence is decided
by a ground condition as a ground sequence. In particular $\rho$ remains
a cardinal and remains regular: a collapse or a shorter cofinal sequence
would supply a new sequence of length less than $\rho$.

The generic union is a new total function $g:\rho\to2$. Totality follows
from the dense sets requiring a value at each coordinate. For every
ground $h:\rho\to2$, the conditions disagreeing with $h$ somewhere are
dense: choose an unused coordinate and assign the opposite value. Hence
$g\notin M$. We have proved
\[
 \delta(M,M[g])=\rho.
\]
Consequently the old full surreal field has a gap of exact cofinality
pair $(\rho,\rho)$ and, for uncountable extension cardinals, is saturated
precisely through $\rho$ and not at $\rho^+$. This realizes every infinite
regular threshold permitted by Proposition~\ref{prop:deltaregular}.
No assertion about preservation of all larger cardinals is needed.

For $\rho=\omega_1^M$, in particular, the old field remains
$\omega_1$-saturated but loses $\omega_2^N$-saturation. There are no new
reals in this case, so $\omega$-saturation survives as well.

\subsection{Prikry forcing separates the two countable notions}

Assume $M$ has a measurable cardinal $\rho$, and use ordinary Prikry
forcing from a normal measure. Its standard properties are that it
preserves cardinals, changes $\cf(\rho)$ to $\omega$, and adds no new
bounded subsets of $\rho$; in particular it adds no reals
\cite{Prikry,Benhamou}. Its generic cofinal $\omega$-sequence in $\rho$
is a new ordinal-valued sequence. Therefore
\[
 \delta(M,N)=\omega,\qquad \R^M=\R^N.
\]
Theorem~\ref{thm:omega} and the main theorem give the strict separation
\[
 K\text{ is }\omega\text{-saturated},\qquad
 K\text{ is not }\omega_1\text{-saturated}.
\]
The omitted countable cut obtained from the interval tree has new
countably many parameters, even though no real is new. The alphabet
$\theta$ in the construction must be allowed to reach $\rho$; restricting
to binary sequences of countable length would miss precisely this example.

At the singularized bound there is another instructive phenomenon:
\[
 (\No_{<\rho})^M=(\No_{<\rho})^N,
\]
because every binary sign of length below $\rho$ codes a bounded subset
of $\rho$. Yet this identical set field is not $\rho$-saturated in $N$:
its all-plus ordinals acquire a countable cofinal sequence. This directly
shows why regularity is an independent clause in
Theorem~\ref{thm:fragments}.

\subsection{Comparison table}

\begin{center}\small
\begin{tabular}{@{}P{0.26\textwidth}P{0.13\textwidth}P{0.21\textwidth}P{0.28\textwidth}@{}}
\toprule
Extension & First length & Old ordered $K$ & Old pure $C=K[i]$\\
\midrule
Adds a Cohen real & $\omega$ & Not $\omega$-saturated & Saturated over every set of parameters\\[4pt]
$\operatorname{Fn}_{<\rho}(\rho,2)$, $\rho>\omega$ regular & $\rho$ & $\kappa$-saturated exactly for uncountable $\kappa\leq\rho$ & Same unrestricted set-saturation\\[4pt]
Prikry at a measurable $\rho$ & $\omega$ & $\omega$-saturated, not $\omega_1$-saturated & Same unrestricted set-saturation\\
\bottomrule
\end{tabular}
\end{center}
Adding conjugation to the last column changes every entry to the
corresponding ordered-field entry. The forcing facts in the table have
been stated with their hypotheses above; the saturation conclusions are
applications of the proved theorems.

\section{Why order gaps do not become Cauchy limits}\label{sec:topology}

The native fine topology of a surreal field uses all positive surreal
radii, not only positive real radii. Its additive uniformity has the
entourages $\abs{x-y}<\varepsilon$ for $\varepsilon\in K_{>0}$.
The internal degeneracy of set-indexed convergence in the full surreal
class is already discussed and formalized in parts of the repository
\cite{RepoFoundations,RepoCatalogue}. The following observation extends
that phenomenon to \emph{new} sets of old numbers.

\begin{proposition}[External discreteness and Cauchy rigidity]\label{prop:topology}
In every pair $M\subseteq N$ under consideration:
\begin{enumerate}[label=\textnormal{(\roman*)}]
\item Every $N$-set $A\subseteq K$ is uniformly discrete and closed
in the fine topology of $K$.
\item Every $N$-set-indexed Cauchy net in $K$ is eventually constant.
Every convergent such net is eventually equal to its limit.
\item The same assertions hold for $C$ with its surreal-valued modulus.
\end{enumerate}
\end{proposition}
\begin{proof}
The nonzero pairwise distances between elements of $A$ form an $N$-set
of positive old numbers. Theorem~\ref{thm:oneside}, with left side
$\{0\}$, gives a positive old lower bound $\varepsilon$ for all of
them. Distinct elements of $A$ are therefore uniformly separated. For
$x\in K\setminus A$, apply the same argument to
$\{\abs{x-a}:a\in A\}$ to obtain a neighborhood of $x$ disjoint from
$A$. This proves closedness.

The range of a set-indexed net is an $N$-set. Choose $\varepsilon>0$
smaller than all nonzero pairwise distances in its range. The Cauchy
condition for this one entourage makes every pair of sufficiently late
values equal. The net is therefore eventually constant. For a convergent
net, apply the separation argument to its range together with its limit.

For $C$, the norm $\abs{a+bi}=\sqrt{a^2+b^2}$ is old for old inputs
and lies in $K_{>0}$ away from zero. Apply the identical argument to
its distance sets.
\end{proof}

For the interval tree, the family
$\{b_{f\res\alpha}-a_{f\res\alpha}:\alpha<\lambda\}$ consists of
positive old numbers and has a common positive old lower bound. The
endpoint sequences are not Cauchy in the fine uniformity. Their omitted
order cut is consequently not evidence of a missing fine-Cauchy limit.
A topology that looks only at selected scales could behave differently,
but it would be a different structure and is not used in the spectrum theorem.

This also explains why elementarity, saturation, and fine completeness
must not be treated as interchangeable notions. The old field remains
real closed and elementary, loses specified types, and still has the
strong set-indexed Cauchy rigidity just proved.

\section{A Boolean completion of reverse simplicity}\label{sec:boolean}

Berenbeim's Question 2 asks for a complete class Boolean algebra in which
the reverse-simplicity order of surreal numbers embeds densely
\cite[p.~32]{Berenbeim}. The answer depends on the meaning of
``complete.'' We give an explicit positive answer for joins of sets,
and a negative answer for joins of all subclasses under Global Choice.

This is a different use of the sign tree from numerical cuts. In this
section, \emph{stronger} means extending a sign, regardless of its numerical
value. Under the identification of $-$ and $+$ with $0$ and $1$, the
partial order is
\[
 P={}^{<\Ord}2,\qquad p\leq_P q\ \Longleftrightarrow\ q\subseteq p.
\]
Each condition is a set; $P$ is a proper class. It is the reverse of the
non-strict initial-segment simplicity order.

\subsection{A direct limit with set-coded elements}

For each ordinal $\alpha$, let
\[
 B_\alpha=\Pow({}^\alpha2).
\]
For $\alpha\leq\beta$, define the cylinder embedding
\begin{equation}\label{eq:cylinder}
 j_{\alpha\beta}(A)=
 \{s\in{}^\beta2:s\res\alpha\in A\}.
\end{equation}
These are injective Boolean homomorphisms, preserve arbitrary set-indexed
unions and intersections, and satisfy
$j_{\beta\gamma}j_{\alpha\beta}=j_{\alpha\gamma}$.

Consider pairs $(\alpha,A)$ with $A\subseteq{}^\alpha2$ and identify
two pairs when their images agree at a common later level. A literal
proper-class equivalence class must not be made an element. Instead,
represent each equivalence class by its unique \emph{least-level pair}.
Such a pair exists: among the levels at most a given representative's
level, there is a least level from which that representative is a cylinder.
Injectivity makes its subset unique. Taking the least-level pair is a
definable operation on the set-coded representatives.

Let $B$ be the class of these canonical pairs. To perform a finite Boolean
operation, lift its arguments to a common level, perform the ordinary
power-set operation there, and take the canonical pair. The result is
independent of the common level by \eqref{eq:cylinder}. Thus $B$ is a
class Boolean algebra whose individual elements are sets.

\begin{theorem}[Set-complete dense completion]\label{thm:booleanyes}
The class Boolean algebra $B$ is complete for every set-sized family
of its elements. The map
\[
 e:P\longrightarrow B\setminus\{0\},\qquad
 e(p)=[(\bd(p),\{p\})],
\]
is an injective dense order embedding. In particular,
\[
 e(p)\leq_B e(q)\quad\Longleftrightarrow\quad q\subseteq p.
\]
\end{theorem}
\begin{proof}
Let $S$ be a set of canonical pairs, and choose an ordinal $\beta$
at least as large as every level occurring in $S$. Lift all of them
to $B_\beta$ and take their union. Its canonical pair is their join
in $B$. Indeed, any alleged upper bound can be moved to a common still
later level, where the ordinary power-set union is below it. Cylinder
embeddings preserve that union, so the comparison remains true in $B$.
Complements give meets as well. The empty join and meet are $0$ and $1$.

A condition $p$ determines its nonempty cylinder of extensions. If $p$
extends $q$, its cylinder is contained in the cylinder of $q$.
If $p$ and $q$ disagree, their cylinders are disjoint. If $p$ is a proper
prefix of $q$, the cylinder of $p$ is not contained in that of $q$:
extend $p$ with the opposite sign at the next prescribed coordinate of
$q$. These cases establish the displayed order equivalence and
injectivity.

Finally, let $b\ne0$ be represented by $(\alpha,A)$ with $A\ne\varnothing$.
Choose $p\in A$. Then $e(p)\leq b$. Thus the embedding is dense.
\end{proof}

The construction answers the stated existence question when ``complete''
means set-complete. This is the size convention in which a class Boolean
algebra has joins for its set-sized subsets, as in the definitions of
Holy--Krapf--L\"ucke--Njegomir--Schlicht \cite{HolyEtAl}. General existence
and uniqueness issues for class forcing are subtle; our construction
works directly for this particular binary tree. We make no claim that
every class partial order admits the same construction, nor that this
completion is the unique class completion.

\subsection{Why joins of all subclasses are impossible}

For the negative statement only, work in GBC with Global Choice.
A class Boolean algebra is \emph{class-complete} if every subclass of
its carrier has a supremum in the carrier, including proper subclasses
in the given class universe. This is strictly stronger than the
set-completeness just proved.

\begin{lemma}[Diagonal obstruction]\label{lem:classdiag}
A class Boolean algebra with an ordinal-indexed proper-class antichain
of nonzero elements cannot be class-complete in GBC with Global Choice.
\end{lemma}
\begin{proof}
Let $(a_\alpha)_{\alpha\in\Ord}$ be such an antichain. The carrier is a
proper class, and Global Choice gives a class surjection
$h:\Ord\to B$; one may use a set-like global well-order and restrict it
to the carrier. Define the class of ordinals
\[
 A=\{\alpha:a_\alpha\not\leq h(\alpha)\}.
\]
If the subclass $\{a_\alpha:\alpha\in A\}$ had a supremum $b$, then
\begin{equation}\label{eq:diagiff}
 a_\alpha\leq b\quad\Longleftrightarrow\quad\alpha\in A.
\end{equation}
The forward implication when $\alpha\notin A$ is ruled out as follows:
$a_\alpha$ is disjoint from every member of the subclass, so
$\neg a_\alpha$ is an upper bound of the entire subclass. Minimality
of its supremum gives $b\leq\neg a_\alpha$. Since $a_\alpha\ne0$,
we cannot also have $a_\alpha\leq b$. The reverse implication in
\eqref{eq:diagiff} follows from the definition of supremum.

Choose $\beta$ with $h(\beta)=b$. The definition of $A$ now gives
\[
 \beta\in A\quad\Longleftrightarrow\quad
 a_\beta\not\leq b\quad\Longleftrightarrow\quad\beta\notin A,
\]
a contradiction. The subclass and all comparisons used here are legitimate
class-comprehension instances; no collection of all classes is formed.
\end{proof}

\begin{theorem}[No all-subclass completion]\label{thm:booleanno}
In GBC with Global Choice there is no class-complete Boolean algebra
with a dense order embedding of $P={}^{<\Ord}2$ into its nonzero part.
\end{theorem}
\begin{proof}
The signs $p_\alpha=+^\alpha{}^\frown-$, for $\alpha\in\Ord$, are
pairwise incompatible. Under a dense order embedding their images are
pairwise disjoint. Indeed, a nonzero intersection would have a condition
below it by density, yielding a common extension of the incompatible
signs. Their images are therefore an ordinal-indexed proper-class
antichain of nonzero Boolean elements. Apply Lemma~\ref{lem:classdiag}.
\end{proof}

The two answers are consistent: a set-sized family of cylinders has a
common set-sized level at which its join can be computed; an arbitrary
proper-class family need not. Set-completeness neither implies nor
purports to imply all-subclass completeness. This distinction also
prevents a misleading claim that Question 2 requires an inconsistency
of ordinary set or class theory.

\section{Scope, significance, and further targets}\label{sec:scope}

\subsection{What the new information consists of}

The usual statement that the full surreal class realizes every set-sized
cut is internal to a universe. The main theorem explains exactly how it
changes when the elements are held fixed while the allowed families of
parameters change. Its most distinctive strengthening is the combination
of three facts:
\begin{enumerate}
\item A cut with one actual ground side is filled even if its other side
is a new and arbitrarily large set.
\item A first new sequence can be encoded by a cut of the same cardinal
length, no matter how large the sequence's ordinal range is.
\item Removing the distinguished real structure by passing to the pure
surcomplex language erases this obstruction entirely.
\end{enumerate}
Together they identify a concrete set-theoretic invariant inside the
external order of the old surreal field. They also demonstrate that
``the model theory of surcomplex numbers'' has no language-independent
saturation answer.

The interval method has a broader range of applicability. It uses only
set-sized interpolation inside the ground order, enough room to split
an interval into any prescribed set of pairwise disjoint subintervals,
and a ground-set containment property for external parameter sets.
Surreal numbers supply all three in an explicit form. It is reasonable
to investigate other definable universal orders with the same properties;
no claim of a classification of all such orders is made here.

\subsection{Limits of the assertions}

The argument does not address exponential or differential expansions of
$\No$. Their types are not classified by real closed cuts, so the
sufficiency proof cannot simply be copied after adding those operations.
Even if their individual old operations were shown absolute, that would
not automatically supply a type-realization theorem in the expanded
language.

The results also do not assert that every class function or arbitrary
proper-class cut is available in ZFC. Set-sized data, a definable ground
predicate, and the indicated uses of class principles are part of the
hypotheses, not conventions to be supplied later. The all-subclass Boolean
obstruction specifically assumes Global Choice; no optimality claim for
that assumption is made.

The article is not a consistency proof for ZFC, GBC, measurable cardinals,
or the existence of the relevant transitive-model configurations.
Theorems are proved under the stated hypotheses; the Prikry example is
conditional on its measurable-cardinal assumption.

No theorem in this manuscript is represented as Lean-checked. The
repository's existing formal proofs and its source manuscripts are
separate evidence. Ordinary written proof is what is supplied here.
The construction of the PDF is reproducible, but successful typesetting
is not formal mathematical verification.

\subsection{A formalization route without class-sized type objects}

A useful first formal target is the branch-decoder theorem. It needs a
set ordinal $\lambda$, a set alphabet $\theta$, a table of intervals,
and a predicate expressing which sequences are old. Abstracting the
closure properties of that predicate isolates the logical core from
surreal arithmetic. The proof obligations are sibling uniqueness,
limit nesting, totality of the default decoder, and its agreement with
a putative externally coded branch.

The next target is the one-old-side theorem. Its main inputs are a
bounded-stage cover for every external set of old elements and an
internal cut operation. The formal theorem should quantify over actual
set-sized sides, not attempt to encode all proper-class cuts as values
of one type.

For a universe-polymorphic proof assistant, the set-fragment theorem of
Section~\ref{sec:fragments} is the cleanest end-to-end theorem: all
structures, types and comparison maps can be represented in a sufficiently
large universe, while the distinction between the two set models is
explicit. The full-class formulations then become metatheorems or
universe-relative schemes rather than an illicit same-universe type of
all surreals.

Finally, the surcomplex result can be separated into two reusable
model-theoretic lemmas: a proper-class algebraically closed field with
ground-set parameter covers has the stated external type realization;
and the quadratic complexification with conjugation interprets the
ordered base field through fixed squares. This separation avoids
repeating the repository's existing quadratic-algebra development.

\subsection{Questions remaining after these results}

One natural next problem is to determine the external saturation spectrum
of old surreal exponential fields. The main theorem supplies a necessary
condition, since the ordered-field reduct must be saturated, but no
sufficiency theorem for that expanded language is proved here.

A second problem is to classify more of the external cut-cofinality
spectrum than its first threshold. The present article proves that no
set-character gap has a side below $\delta$ and constructs one at
$(\delta,\delta)$. It does not characterize all pairs above that level
or decide which asymmetrical pairs can occur for a specified forcing.

A third problem is to identify weak class-theoretic assumptions under
which one can construct an isomorphism of the two pure surcomplex class
fields $C$ and $C^+$. The set-saturation theorem provides local
realizations but is not itself a global class back-and-forth construction.
Each of these questions remains distinct from the theorems proved here.

\appendix
\section{Repository and literature comparison}\label{app:audit}

\subsection{Pinned source review}

The repository review used the snapshot
\[
 \fullhash
\]
of \path{VladimirReshetnikov/Surreal}. The repository was read through
the connected GitHub interface. It was not modified or independently
built. The recursive tree response was too large to be displayed in full;
it is not treated as evidence of a complete file-by-file audit.

The file reads were pinned to this snapshot; the connector's code searches
used its default-branch index rather than a commit-specific full-text index.
The following records describe the useful scope of the inspection, rather
than asserting that every theorem in every report was searched exhaustively.

\begin{longtable}{@{}P{0.36\textwidth}P{0.56\textwidth}@{}}
\toprule
Source or search & Scope relevant to this article\\
\midrule
\endfirsthead
\toprule
Source or search & Scope relevant to this article\\
\midrule
\endhead
\path{docs/README.md} & The report catalogue and reading routes. It lists forty catalogued reports plus four newly placed drafts, separating foundational, analytic, algebraic and computational directions.\\[5pt]
\path{docs/foundations-and-computation/foundations/article.tex} & Opening 180 source lines, including the abstract and scope declarations. These cover sets, classes, universes, support and topology, rather than the external sequence spectrum proved here. This is not a claim that every later section was read.\\[5pt]
\path{docs/surcomplex/surcomplex-field-automorphisms/article.tex} & The returned opening text: abstract, structural distinctions, real-axis stabilizer and initial conjugation discussion. The real-axis/pure-field distinction is acknowledged as existing material, not repackaged as a new observation.\\[5pt]
Repository root \path{README.md} & Returned foundational and algebraic implementation overview. Used to avoid implying that this manuscript supplied the existing sign or quadratic-field infrastructure.\\[5pt]
Scoped code searches & Searches for \texttt{forcing}, \texttt{saturation} and \texttt{distributivity} returned no matches. These are limited search results, not a proof that the words or equivalent results are absent from every file or past revision.\\
\bottomrule
\end{longtable}

The proposed contribution is therefore specifically the new sequence/gap
spectrum and its forcing-dependent language comparison, not the elementary
facts that real closed fields have an algebraically closed quadratic
extension or that conjugation singles out a real axis. The topological
section explicitly starts from an internally familiar repository phenomenon
and proves its persistence for new sets of old values.

\subsection{Literature boundary}

Conway and Gonshor supply the basic construction and cut theory.
Van den Dries--Ehrlich, including their erratum, and
Bournez--Guilmant supply the birthday-field background.
Marker supplies the standard quantifier-elimination and saturation
machinery. Jech supplies the forcing preservation vocabulary.
Prikry and Benhamou supply the large-cardinal forcing example.
Berenbeim supplies the explicit reverse-simplicity Boolean question and
the earlier suggestion to study surreals under forcing.
Holy and coauthors supply the surrounding class-forcing completion theory.

Targeted web searches for the combination of surreal numbers, forcing,
ground models and saturation did not locate a prior statement of the
exact theorems proved here. Search results were not uniformly relevant,
and the investigation is not an exhaustive mathematical bibliography.
Accordingly, the appropriate novelty claim is:
\begin{quote}
These are proved candidate contributions not located in the consulted
literature or the stated repository review, subject to independent
mathematical checking and further priority research.
\end{quote}
That qualification concerns verification and priority. It does not replace
any hypothesis or proof step in the theorems.

\section{Dependency and proof-obligation ledger}\label{app:ledger}

\begin{longtable}{@{}P{0.27\textwidth}P{0.34\textwidth}P{0.29\textwidth}@{}}
\toprule
Result & Essential inputs & Boundary checked\\
\midrule
\endfirsthead
\toprule
Result & Essential inputs & Boundary checked\\
\midrule
\endhead
Old arithmetic and elementarity & Absolute signs and set-like Conway recursion; real closed and algebraically closed QE & Field elementarity is not universe elementarity.\\[5pt]
Ground-set containment & Birthday function; common ordinals; old bounded sign sets & The containing ground set need not have the size of the external set.\\[5pt]
One old side & Ground-set containment; ground cut theorem & Both sides are external sets; one is literally in $M$.\\[5pt]
Interval tree and decoder & Set-length recursion; explicit successor splitting; two cuts at each limit node & Decoder has a default value and is a total function in $M$.\\[5pt]
Main spectrum & Minimal new sequence; decoder; one-side lemma; RCF QE & The $\eta_\kappa$/saturation equivalence is used only for uncountable $\kappa$.\\[5pt]
First gap & Minimal-length regularity; nested endpoints; one-side interpolation & Exact cofinalities are computed in $N$ and have set witnesses.\\[5pt]
$\omega$-saturation & Countable syntax over finite old parameters; rational cuts & No new reals is not replaced by no new countable ordinal sequences.\\[5pt]
Pure surcomplex saturation & Ground-set containment; old algebraic closures of set fields; ACF QE & No class-sized transcendence basis or global isomorphism is assumed.\\[5pt]
Conjugation spectrum & Fixed field; order by nonzero squares; coordinate interpretation & Naming $i$ adds only one parameter; finite-variable types are used.\\[5pt]
Bounded fragments & Classical real closedness; birthday cut bound; sign-cone lemma & Regularity in $N$ is required separately from equality of short signs.\\[5pt]
External Cauchy rigidity & One-old-side bound on all positive distances & Fine surreal radii, not only real or countably chosen radii.\\[5pt]
Set-complete Boolean algebra & Power-set cylinder maps; bounded levels for set families & Quotients are represented by least-level set pairs, not class objects.\\[5pt]
No class-complete algebra & Global Choice; proper-class antichain; diagonal subclass & This is a stronger completeness requirement, with an extra stated axiom.\\
\bottomrule
\end{longtable}

The ledger records the mathematical checks performed in assembling the
arguments. It is not a proof-assistant log. In particular, finite numerical
experiments would not verify the transfinite decoder or the class-size
assertions, so none are presented as substitutes for their proofs.

\phantomsection
\addcontentsline{toc}{section}{References}
\begin{thebibliography}{99}

\bibitem{Conway}
J. H. Conway, \emph{On Numbers and Games}, 2nd ed., A K Peters, 2001
(first edition, Academic Press, 1976).

\bibitem{Gonshor}
H. Gonshor, \emph{An Introduction to the Theory of Surreal Numbers},
London Mathematical Society Lecture Note Series 110,
Cambridge University Press, 1986.

\bibitem{vdDE}
L. van den Dries and P. Ehrlich,
``Fields of surreal numbers and exponentiation,''
\emph{Fundamenta Mathematicae} \textbf{167} (2001), 173--188.
\url{https://www.impan.pl/shop/en/publication/transaction/download/product/89226}.

\bibitem{vdDEerr}
L. van den Dries and P. Ehrlich,
``Erratum to `Fields of surreal numbers and exponentiation',''
\emph{Fundamenta Mathematicae} \textbf{168} (2001), 295--297.
The corrected background statements used here are also restated in
\cite[Theorems 1.1--1.2 and 3.3]{BournezGuilmant}.

\bibitem{BournezGuilmant}
O. Bournez and Q. Guilmant,
``Surreal fields stable under exponential and logarithmic functions,''
arXiv:2201.08199, 2022.
\url{https://arxiv.org/abs/2201.08199}.

\bibitem{Marker}
D. Marker, \emph{Model Theory: An Introduction},
Graduate Texts in Mathematics 217, Springer, 2002.
\url{https://doi.org/10.1007/b98860}.

\bibitem{Jech}
T. Jech, \emph{Set Theory: The Third Millennium Edition, Revised and
Expanded}, Springer Monographs in Mathematics, Springer, 2003.

\bibitem{Prikry}
K. L. Prikry, ``Changing measurable into accessible cardinals,''
\emph{Dissertationes Mathematicae (Rozprawy Matematyczne)}
\textbf{68} (1970), 55 pp.

\bibitem{Benhamou}
T. Benhamou, ``Prikry forcing and tree Prikry forcing of various filters,''
arXiv:1801.04424, 2018; author-hosted published version, 2019.
\url{https://arxiv.org/abs/1801.04424};
\url{https://sites.math.rutgers.edu/~tb822/Prikry_forcing_and_Tree_Prikry.pdf}.

\bibitem{Berenbeim}
A. Berenbeim, \emph{Review of Surreal Number Construction and Forcing
Construction}, research notes dated 8 April 2020, 35 pp.
Question 2, p.~32; proposed further directions, p.~34.
Listed by the author as incomplete notes.
\url{https://www.berenbeim.com/papers};
\url{https://www.berenbeim.com/_files/ugd/4185a9_2da51607fad1430fa1735144805340dd.pdf}.

\bibitem{HolyEtAl}
P. Holy, R. Krapf, P. L\"ucke, A. Njegomir and P. Schlicht,
``Class forcing, the forcing theorem and Boolean completions,''
arXiv:1710.10820, 2017.
\url{https://arxiv.org/abs/1710.10820}.
In particular, see the definitions of model-completeness and Boolean
completion in the introduction; the paper also treats nonuniqueness.

\bibitem{RepoCatalogue}
V. Reshetnikov and contributors, \emph{Surreal} repository,
\path{docs/README.md}, snapshot \shorthash.
\url{https://github.com/VladimirReshetnikov/Surreal/blob/4f2645599121fa872c7104995e47f86f0382351f/docs/README.md}.

\bibitem{RepoFoundations}
\emph{Foundations for Surreal and Surcomplex Mathematics},
AI-assisted research exposition in the \emph{Surreal} repository,
\path{docs/foundations-and-computation/foundations/article.tex},
snapshot \shorthash. The review scope is stated in
Appendix~\ref{app:audit}.

\bibitem{RepoAutomorphisms}
\emph{Automorphisms of the Surcomplex Numbers},
AI-assisted research exposition in the \emph{Surreal} repository,
\path{docs/surcomplex/surcomplex-field-automorphisms/article.tex},
snapshot \shorthash. The review scope is stated in
Appendix~\ref{app:audit}.

\end{thebibliography}
\end{document}
